% 图 53.4 —— 由 `checks/emit_hand.py` 自动拼出（probe 精测 + prims2 图元分解）
%%
% ⚠ 这是**稿子不是成品**：跑 `checks/checkhand.py 53.4` 看指标 + 看叠合图。
%%
%% ⚠ 待核：
%   · 本图 probe 报出阴影族 1 个 —— **本脚本不生成阴影**（要照原书手写）；prims2 的 strokes 里可能含阴影线，核对时留意别画重
%   · 可疑字形 2 项（空心圈 ⊙ 常被认成字母 o）—— 见文件末
%%
\documentclass[12pt]{standalone}
\usepackage{fontspec}
\setsansfont{Arial}
\newfontfamily\mfallback{DejaVu Sans}
\usepackage{tikz}
\begin{document}
\begin{tikzpicture}[x=1pt, y=-1pt, line width=4.0pt, line cap=round, line join=round]

  %% ════ 填充区（prims2 的 fills）════

  %% ════ 直线段（prims2 的 strokes）════
  \draw[line width=4.0pt] (466.0,108.0) -- (551.0,108.0);
  \draw[line width=5.0pt] (145.0,269.0) -- (145.0,180.0);
  \draw[line width=5.0pt] (795.0,227.0) -- (793.0,232.0);
  \draw[line width=4.0pt] (55.0,270.0) -- (56.0,180.0);
  \draw[line width=5.0pt] (238.0,270.0) -- (238.0,180.0);
  \draw[line width=4.0pt] (328.0,86.0) -- (349.0,85.0);
  \draw[line width=5.0pt] (327.0,87.0) -- (327.0,177.0);
  \draw[line width=5.0pt] (552.0,107.0) -- (548.0,101.0);
  \draw[line width=4.0pt] (239.0,388.0) -- (238.0,363.0);
  \draw[line width=5.0pt] (239.0,271.0) -- (325.0,270.0);
  \draw[line width=4.0pt] (238.0,272.0) -- (238.0,361.0);
  \draw[line width=6.0pt] (552.0,109.0) -- (549.0,114.0);
  \draw[line width=5.0pt] (37.0,179.0) -- (55.0,179.0);
  \draw[line width=5.0pt] (146.0,87.0) -- (238.0,86.0);
  \draw[line width=5.0pt] (793.0,157.0) -- (776.0,156.0);
  \draw[line width=4.0pt] (239.0,88.0) -- (238.0,178.0);
  \draw[line width=5.0pt] (804.0,195.0) -- (803.0,189.0);
  \draw[line width=5.0pt] (803.0,196.0) -- (761.0,193.0);
  \draw[line width=6.0pt] (571.0,108.0) -- (553.0,108.0);
  \draw[line width=5.0pt] (240.0,87.0) -- (326.0,86.0);
  \draw[line width=5.0pt] (416.0,270.0) -- (328.0,271.0);
  \draw[line width=4.0pt] (325.0,362.0) -- (239.0,362.0);
  \draw[line width=4.0pt] (417.0,271.0) -- (417.0,327.0);
  \draw[line width=5.0pt] (950.0,141.0) -- (953.0,128.2);
  \draw[line width=4.0pt] (953.0,128.2) -- (946.8,103.8);
  \draw[line width=4.0pt] (946.8,103.8) -- (928.1,85.1);
  \draw[line width=5.0pt] (678.0,91.0) -- (661.0,115.0);
  \draw[line width=4.0pt] (146.0,66.0) -- (145.0,85.0);
  \draw[line width=4.0pt] (326.0,363.0) -- (327.0,412.0);
  \draw[line width=5.0pt] (327.0,43.0) -- (327.0,85.0);
  \draw[line width=3.0pt] (123.0,87.0) -- (144.0,86.0);
  \draw[line width=4.0pt] (326.0,178.0) -- (239.0,179.0);
  \draw[line width=4.0pt] (418.0,270.0) -- (475.0,269.0);
  \draw[line width=4.0pt] (56.0,361.0) -- (56.0,272.0);
  \draw[line width=5.0pt] (145.0,88.0) -- (145.0,177.0);
  \draw[line width=5.0pt] (328.0,178.0) -- (416.0,178.0);
  \draw[line width=4.0pt] (327.0,179.0) -- (326.0,269.0);
  \draw[line width=5.0pt] (56.0,178.0) -- (57.0,167.0);
  \draw[line width=4.0pt] (622.0,144.3) -- (625.0,114.5);
  \draw[line width=4.0pt] (625.0,114.5) -- (640.2,87.8);
  \draw[line width=5.0pt] (57.0,362.0) -- (144.0,363.0);
  \draw[line width=5.0pt] (146.0,179.0) -- (237.0,179.0);
  \draw[line width=5.0pt] (239.0,85.0) -- (238.0,16.0);
  \draw[line width=5.0pt] (57.0,179.0) -- (144.0,178.0);
  \draw[line width=5.0pt] (760.0,191.0) -- (775.0,157.0);
  \draw[line width=5.0pt] (803.0,208.0) -- (799.0,217.0);
  \draw[line width=5.0pt] (56.0,381.0) -- (55.0,363.0);
  \draw[line width=4.0pt] (16.0,271.0) -- (54.0,271.0);
  \draw[line width=5.0pt] (327.0,271.0) -- (327.0,361.0);
  \draw[line width=4.0pt] (417.0,140.0) -- (417.0,177.0);
  \draw[line width=4.0pt] (36.0,362.0) -- (54.0,362.0);
  \draw[line width=4.0pt] (145.0,415.0) -- (145.0,364.0);
  \draw[line width=5.0pt] (237.0,362.0) -- (146.0,363.0);
  \draw[line width=5.0pt] (146.0,271.0) -- (237.0,271.0);
  \draw[line width=5.0pt] (145.0,272.0) -- (145.0,362.0);
  \draw[line width=4.0pt] (328.0,362.0) -- (355.0,362.0);
  \draw[line width=5.0pt] (57.0,271.0) -- (144.0,270.0);
  \draw[line width=5.0pt] (457.0,178.0) -- (418.0,178.0);
  \draw[line width=5.0pt] (417.0,179.0) -- (417.0,269.0);

  %% ════ 曲线（prims2 的 curves —— 三段式，坐标同为 (y,x)）════
  \draw[line width=5.0pt] (712.0,124.0) .. controls (726.5,145.5) and (752.9,150.0) .. (775.0,156.0);
  \draw[line width=5.0pt] (734.0,134.0) .. controls (766.6,91.9) and (841.8,90.3) .. (876.0,131.0);
  \draw[line width=5.0pt] (795.0,157.0) .. controls (833.9,155.7) and (877.6,153.7) .. (902.0,120.0);
  \draw[line width=5.0pt] (805.0,196.0) .. controls (858.6,195.8) and (918.3,188.3) .. (950.0,141.0);
  \draw[line width=5.0pt] (928.1,85.1) .. controls (855.5,40.6) and (748.5,42.5) .. (678.0,91.0);
  \draw[line width=4.0pt] (661.0,115.0) .. controls (658.1,125.3) and (660.7,136.3) .. (666.1,145.1);
  \draw[line width=4.0pt] (666.1,145.1) .. controls (688.2,174.1) and (724.9,185.7) .. (759.0,191.0);
  \draw[line width=4.0pt] (769.0,240.0) .. controls (762.6,226.0) and (756.1,210.2) .. (760.0,194.0);
  \draw[line width=4.0pt] (768.0,241.0) .. controls (710.1,233.5) and (637.9,208.5) .. (622.0,144.3);
  \draw[line width=4.0pt] (640.2,87.8) .. controls (699.0,25.7) and (790.7,18.2) .. (871.4,32.0);
  \draw[line width=4.0pt] (871.4,32.0) .. controls (918.8,41.9) and (965.5,69.7) .. (993.0,113.0);
  \draw[line width=5.0pt] (993.0,113.0) .. controls (997.3,127.9) and (997.2,145.6) .. (991.9,160.1);
  \draw[line width=5.0pt] (991.9,160.1) .. controls (943.0,233.0) and (852.3,252.4) .. (770.0,241.0);

  %% ════ 实心点 ════
  \fill (794.5,162.0) circle (3.6);

  %% ════ 标签（probe 直接给的 \node 行）════
  %% ⚠ 全部补了 fill=white：文字在最上层，否则与曲线交叉干扰
     \node[anchor=center,inner sep=0pt,fill=white,font=\fontsize{33.8}{42.3}\selectfont] at (510.1,81.5) {\textsf{\textbf{×}}};   % box(502,71,15,16)
     \node[anchor=center,inner sep=0pt,fill=white,font=\fontsize{29.2}{36.4}\selectfont] at (481.9,84.0) {\textsf{\textbf{\textit{P}}}};   % box(471,72,18,22)
     \node[anchor=center,inner sep=0pt,fill=white,font=\fontsize{29.2}{36.4}\selectfont] at (541.9,85.0) {\textsf{\textbf{\textit{P}}}};   % box(531,73,18,22)
     \node[anchor=center,inner sep=0pt,fill=white,font=\fontsize{42.8}{53.5}\selectfont] at (801.8,188.7) {\textsf{\textbf{′}}};   % box(799,172,7,12)
     \node[anchor=center,inner sep=0pt,fill=white,font=\fontsize{25.7}{32.1}\selectfont] at (427.1,355.9) {\textsf{\textbf{2}}};   % box(420,347,13,17)
     \node[anchor=center,inner sep=0pt,fill=white,font=\fontsize{30.0}{37.5}\selectfont] at (408.0,364.0) {\textsf{\textbf{\textit{R}}}};   % box(398,352,20,22)

  % 钉死画布：否则超出的图元/控制点会把页面撑大，1:1 回比就失准
  \pgfresetboundingbox
  \path[use as bounding box] (0,0) rectangle (1011,427);
\end{tikzpicture}
\end{document}

% ⚠ 待核清单（probe 拿不准，人眼过一遍）：
%   · 未识别/跳过：⑤ 跳过/未识别的字形
%   · 未识别/跳过：box(790,225,9,10)
